\documentclass[%
 reprint,
 amsmath,amssymb,
 aps,
]{revtex4-1}

\usepackage{graphicx}% Include figure files
\usepackage{dcolumn}% Align table columns on decimal point
\usepackage{bm}% bold math
\usepackage[dvipdfm,colorlinks=true,linkcolor=blue,filecolor=blue, urlcolor=blue]{hyperref}
\usepackage{braket}
\usepackage{caption} % for captionsetup[figure]
\usepackage{threeparttable} % for table caption modification
\usepackage{breqn} % dmath: automatic equation line breaking
\numberwithin{equation}{section}% (p.37)

% \sloppy sets all text to loose spacing, allowing overflow and uneven word spacing;
% not recommended for academic papers requiring strict typesetting.
\tolerance=1000
\emergencystretch=1em
\hyphenpenalty=3000
\exhyphenpenalty=3000
\flushbottom
\usepackage[final]{microtype}


\begin{document}

    \captionsetup[figure]{justification=raggedright,
        labelfont={bf},name={Fig.},labelsep=period} % raggedright for left-aligned figure captions
    \captionsetup[table]{
        labelfont={bf},name={Table.},labelsep=period}


\title{From Clock Synchronization to Electromagnetism: \\A Realist Construction of $U(1)$ Gauge Theory}

\author{Satoshi Hanamura}
\email{hana.tensor@gmail.com}
\date{June 1, 2025}



\begin{abstract}
Gauge symmetry is a cornerstone of modern field theory, yet its conceptual status remains unsettled: conventional formulations treat it as a redundancy in mathematical representation rather than a physically grounded structure. Here we propose a realist reinterpretation of $U(1)$ gauge symmetry in electromagnetism, demonstrating how gauge freedom emerges from internal phase dynamics intrinsic to the electron. Within the 0-Sphere model, the electron is modeled not as a point particle but as a temporally structured system possessing internal harmonic motion between two thermal potential energy kernels. This internal structure gives rise to deterministic, $4\pi$-periodic dynamics and a Zitterbewegung velocity $v_{\text{ZB}} = 0.040374c$, derived from first principles by combining special relativity with geodesic precession from general relativity.
Coherent phase transport across an ensemble of such particles naturally induces a $U(1)$ gauge connection, interpreted as a geometric synchronization mechanism rather than a formal constraint. Furthermore, the internal energy structure encodes both $\omega t$ and $\omega t/2$ components, establishing a harmonic hierarchy that bridges bosonic gauge behavior with fermionic spin. This framework offers a particle-based, ontologically grounded reinterpretation of gauge symmetry, in which electromagnetic interactions emerge from the holonomy of internal phase transport. The resulting picture opens a path toward unifying quantum structure with geometric principles of interaction, without invoking non-Abelian complexity.
\end{abstract}


\maketitle

% ========================================
% Section 1: Introduction
% ========================================
\section{Introduction}

Fundamental interactions in modern physics are described by gauge theories, yet their gauge symmetry remains conceptually problematic. Gauge transformations connect mathematically equivalent representations of the same physical state, introducing redundancy without clear physical significance~\cite{zee2010}. In quantum electrodynamics, the $U(1)$ phase freedom of the wavefunction necessitates gauge potentials, but their microscopic origin is often imposed axiomatically rather than derived, leaving their physical role ambiguous.

Various approaches have sought to address this conceptual redundancy. Lattice gauge theory eliminates local gauge freedom by formulating interactions on discrete spacetime points~\cite{wilson1974}, while loop quantum gravity re-expresses gauge fields through non-redundant, holonomy-based variables such as Wilson loops~\cite{rovelli1995}. These frameworks offer structural solutions but rely on discretization or background independence. In contrast, we propose a continuous, particle-centric model in which the $U(1)$ gauge connection emerges from synchronized internal phase dynamics.


We propose a realist reinterpretation of $U(1)$ gauge symmetry, grounded in the internal phase dynamics of electrons within the 0-Sphere model~\cite{hanamura2024quantum}. Electrons, modeled as oscillators with thermal potential energy (TPE) kernels---localized energy reservoirs---generate a geometric gauge connection through synchronized oscillations. The Aharonov-Bohm effect demonstrates the physical significance of gauge potentials~\cite{aharonov1959}, supporting our particle-based approach. Inspired by foundational work in gauge theory~\cite{weyl1929, yang1954}, this framework unifies quantum mechanics with electromagnetic interactions.



% ========================================
% Section 2: Motivation and Conceptual Tension in Gauge Theory
% ========================================
\section{Motivation and Conceptual Tension in Gauge Theory}

In the conventional interpretation of gauge theory, the gauge transformation of the wavefunction,
\begin{equation}
\psi(x) \rightarrow e^{i\alpha(x)} \psi(x),
\label{eq:wavefunction}
\end{equation}
is accompanied by a transformation of the gauge potential, 
\begin{equation}
A_\mu(x) \rightarrow A_\mu(x) - \frac{1}{e} \partial_\mu \alpha(x),
\label{eq:gauge_potential}
\end{equation}
ensuring the invariance of the covariant derivative $D_\mu \psi = (\partial_\mu + i e A_\mu) \psi$. However, this structure introduces a paradox: the gauge field appears as a formal necessity, with no intrinsic dynamics apart from its role in maintaining local phase invariance. The Aharonov-Bohm effect provided the first compelling evidence that gauge potentials possess physical significance beyond their classical role as mathematical conveniences. Tonomura's definitive experimental demonstration~\cite{tonomura1986} showed that electrons can be influenced by electromagnetic potentials even in regions where classical fields vanish, suggesting that the gauge structure encodes genuine physical information rather than mere calculational redundancy. The connection $A_\mu$ is introduced not because of a physical mechanism, but because of a demand for local symmetry in the mathematical formalism.

Gauge symmetry presents a fundamental conceptual puzzle. While mathematically 
indispensable, it appears to describe not physical reality but mathematical 
redundancy. As Zee observes~\cite{zee2010}:
\begin{quote}
\textit{``But perhaps the most unsatisfying feature of field theory is the present formulation of gauge theories. Gauge ``symmetry'' does not relate two different physical states, but two descriptions of the same physical state. We have this strange language full of redundancy we can't live without. We start with unneeded baggage that we then gauge-fix away. We even know how to avoid this redundancy from the start but at the price of discretizing spacetime.''}
\end{quote}

This redundancy becomes particularly problematic when we consider what gauge 
transformations actually represent physically. This concern has been echoed in the philosophical literature as well. Teller~\cite{teller2000}, for example, argues that gauge symmetry does not reflect a physical transformation between distinct states but rather encodes descriptive redundancy, raising questions about its ontological status. In response to this, we consider a particle-centric model in which the gauge degree of freedom arises from an internal physical mechanism.

In the 0-Sphere model~\cite{hanamura2018}, the electron possesses two thermal potential energy (TPE) kernels that exchange mass-energy in a periodic fashion. This internal motion defines a time-dependent phase $\theta(t) = \omega t$, leading to an internal energy profile given by~\cite{hanamura2024quantum}:
\begin{equation}
E(t) = E_0 \left[ \cos^4\left( \frac{\omega t}{2} \right) + \sin^4\left( \frac{\omega t}{2} \right) + \frac{1}{2} \sin^2(\omega t) \right],
\label{eq:internal_energy}
\end{equation}
which is strictly conserved and exhibits $4\pi$ periodicity. This structure offers a natural explanation for fermionic double-valuedness and the emergence of observable quantities such as the Zitterbewegung.

Moreover, when energy transfer between these TPE kernels is analyzed geometrically, it obeys a Snell-like refraction condition~\cite{hanamura2024quantum},
\begin{equation}
\frac{\sin \theta_A}{\sin \theta_B} = \frac{v_A}{v_B},
\label{eq:snell}
\end{equation}
where the internal phase velocities $v_A$ and $v_B$ are determined by local thermal gradients. This principle, in turn, leads to geodesic phase transport governed by a connection—identified with the electromagnetic gauge potential.

Thus, instead of viewing gauge symmetry as a formal redundancy, we propose it reflects a real physical structure: the geometry of internal phase synchronization. This reinterpretation not only grounds the $U(1)$ gauge field in particle dynamics, but also points toward a realist, experimentally accessible foundation for electromagnetic interaction.

In classical mechanics, the motion of a point particle is fully determined by its position and momentum as functions of time. These quantities form a curve in the tangent bundle $TM$ over the configuration space  $M$, with position corresponding to a point in the base manifold and momentum to a vector in the fiber. This geometric formulation reveals how physical degrees of freedom are encoded in the structure of fiber bundles, where the choice of fiber reflects the nature of the quantity—such as position, velocity, spin, or charge.

This geometric picture extends to field theory, where spacetime becomes the base manifold and internal degrees of freedom form the fiber. In electromagnetism, gauge transformations correspond to changes in the local trivialization of the fiber bundle. Traditionally, these were viewed as redundancies, lacking direct physical consequences.

However, the Aharonov-Bohm effect~\cite{aharonov1959} suggests that this redundancy may harbor genuine physical content. This insight was later extended by Berry~\cite{berry1984}, who showed that geometric phase factors can accumulate in adiabatically evolving quantum systems, thereby linking gauge structure to observable holonomies in parameter space. Tonomura's experiments~\cite{tonomura1986} demonstrated that gauge potentials influence quantum interference even where classical fields vanish, hinting that gauge structure encodes real physical information rather than mere mathematical convenience. This experimental evidence motivates us to seek a physical foundation for gauge freedom.

Quantum gauge theory, by contrast, treats gauge symmetry as fundamental. Local phase invariance underlies charge conservation and necessitates gauge fields with real dynamical significance. Yet this conceptual leap—from redundancy in classical potentials to a physically meaningful symmetry in quantum amplitudes---has not been fully justified. It remains unclear whether gauge transformations are merely computational artifacts inherited from classical formalism or instead reflect intrinsic physical structures.

The conventional framework circumvents this tension by compartmentalizing the classical and quantum domains, unifying them only through the correspondence principle. However, such a division may reflect not a fundamental dichotomy in nature, but a lack of physical grounding in the underlying mathematical structures.

The present work proposes an alternative: that gauge freedom originates in the internal temporal dynamics of particles themselves. In the 0-Sphere model, each particle is endowed with an internal phase structure evolving in proper time, analogous to a clock. This phase is not imposed externally, but arises from a well-defined deterministic oscillation intrinsic to the particle’s ontology. Gauge transformations then acquire a concrete physical interpretation: they represent synchronizations or shifts in this internal clock phase, akin to redefining the origin of time for a localized periodic process.

This interpretation has a clear geometric analog. In classical mechanics, a particle’s motion defines a trajectory in the tangent bundle over spacetime, where each point in the base manifold corresponds to a position, and the associated tangent vector encodes its velocity. Similarly, the internal phase evolution defines a trajectory in a fiber bundle whose base is spacetime and whose fiber encodes internal degrees of freedom—here, the proper-time phase associated with the particle’s internal oscillation, which behaves analogously to an internal clock. Just as in classical mechanics, the choice of local inertial frame in a tangent bundle reflects the freedom to redefine velocity coordinates without altering physical predictions, the gauge degree of freedom in our framework reflects the freedom to redefine internal phase origins across spacetime --- while preserving consistency of physical law as expressed through covariant synchronization.

From this perspective, gauge symmetry ceases to be a purely formal redundancy and instead emerges as an expression of physical structure—one that unifies the computational utility of classical gauge choice with the geometric necessity of quantum field theory. This realist interpretation not only bridges the conceptual gap between classical and quantum gauge theories, but also opens the door to interpreting other fundamental symmetries as manifestations of internal geometric order within matter itself.


% ========================================
% Section 3: Internal Structure and Dynamics of the Electron
% ========================================
\section{Internal Structure and Dynamics of the Electron}

In the 0-Sphere model, the electron is not treated as a point particle but as an oscillator arising from the periodic exchange of energy between two distinct components: a radiation kernel and an absorption kernel. These are identified with the positive- and negative-energy solutions of the Dirac equation, reinterpreted as real-space structures in continuous time.

The internal dynamics is governed by a deterministic phase variable $\theta(t)$, which drives the oscillatory motion between the two kernels. This harmonic-like motion leads to an effective internal velocity, calculated as $v_{\mathrm{ZB}} = 0.040374c$. Although well below the speed of light, this value emerges from the model’s closed Hamiltonian dynamics, incorporating relativistic effects such as internal-state precession and a geodesic precession analysis applied to the kernel’s critical radius.

This framework provides a self-contained dynamical origin for the electron’s intrinsic periodicity, independent of quantum field theoretic assumptions. It offers a physical interpretation of Zitterbewegung---rapid oscillations originally predicted by Schrödinger~\cite{zitterbewegung_schrodinger1930}---as a real oscillatory motion between localized substructures, mediated by the internal clock $\theta(t)$, rather than as a mere artifact of interference in the Dirac formalism.

The 0-Sphere model describes the electron not as a point particle or an extended field excitation, but as a bound interference pattern between a positive-energy radiative kernel and a negative-energy absorptive kernel. These two components oscillate with a phase difference of $\pi / 2$ and form a standing wave defined along the geodesic connecting two discrete points—representing the antipodes of a topological 0-sphere. The kernels are separated by a distance on the order of the Compton wavelength and are dynamically coupled through a surrounding photon sphere. This internal structure naturally gives rise to Zitterbewegung but here interpreted as genuine physical motion between discrete energy kernels, tied together through the internal clock of $\theta (t)$. Such a model provides a self-contained dynamical basis for the electron's internal periodicity, without invoking field-theoretic degrees of freedom.

The finite size of each electron kernel, derived from relativistic and gravitational constraints in our model, provides a natural length scale of approximately \(3.43 \times 10^{-25} \, \text{m}\)~\cite{hanamura2024quantum}---much smaller than the Compton wavelength but finite nonetheless. This prediction is obtained by combining Lorentz contraction from special relativity with geodetic precession from general relativity, linking internal structure to both inertial and gravitational principles. The internal temporal phase $\omega t$ governs oscillatory energy exchange between the two kernels, endowing the electron with a physically meaningful internal clock.

The relative phase between the two kernels evolves in time as a deterministic oscillator. This gives rise to a physically meaningful internal clock, or time phase $\theta(t)$, associated with the particle’s state. The evolution of this phase is governed by a closed-form Hamiltonian previously derived~\cite{hanamura2023}, and leads directly to the observable phenomenon of Zitterbewegung.

This internal phase introduces a new dimension to the physical description of particles. While conventional gauge theories model interactions via field-mediated boson exchange across spacetime, the present approach assigns to each particle a self-contained internal evolution governed by its own phase dynamics. In this framework, interaction may be understood not as a process of continuous field propagation, but as the synchronization of internal time phases across particles.

Although the current formulation focuses on the internal structure of a single electron, it suggests a new mechanism for interaction: the photon sphere that mediates kinetic energy around the TPE kernel can, under certain conditions, detach and transfer to another electron. This dynamic exchange is treated as a physically real process, not merely a virtual field fluctuation. Moreover, the number of photons absorbed or emitted by the photon sphere is not constrained by chemical potential conservation in the usual sense, reflecting a non-conservative interaction framework embedded in the internal oscillatory geometry. Such a structure may provide a future foundation for deriving interparticle gauge interactions from purely particle-based principles.

Accordingly, in this model, interparticle interaction is mediated by the real emission and absorption of photons that detach from the photon sphere surrounding a TPE kernel and deliver kinetic energy to neighboring electrons. In this regard, the picture follows the structure of Feynman diagrams, in which gauge bosons represent physically exchanged quanta between fermionic lines. However, unlike in conventional field theory, these photons are not virtual constructs arising from propagators, but real, geometrically localized energy carriers. This physical picture renders the invocation of a continuous thermal potential energy field unnecessary. The electromagnetic coupling emerges not from an underlying continuum, but from discrete, geometrically constrained transfers of energy between internal oscillators.

This reinterpretation leads to the hypothesis that $U(1)$ gauge symmetry reflects the freedom of choosing the phase origin $\theta \to \theta + \alpha(x)$ of each particle's internal clock. Crucially, unlike traditional formulations where such transformations are considered gauge redundancies, our model treats them as physically real: different phase choices correspond to different initial conditions in the particle’s internal evolution, and the gauge potential arises as a geometric mechanism to coordinate them.

The total energy of the system remains exactly conserved. Let \( E_0 \) denote the rest energy of the electron, i.e., \( E_0 = mc^2 \). Under this assumption, the following identity holds~\cite{hanamura2024quantum, hanamura2018}:
\begin{equation}
E_0 = E_0\left[\cos^4\left(\frac{\omega t}{2}\right) + \sin^4\left(\frac{\omega t}{2}\right) + \frac{1}{2}\sin^2(\omega t)\right].
\label{eq:energy_conservation}
\end{equation}

\textbf{This expression guarantees exact conservation of total energy without invoking external fields or continuous media.} At the same time, it introduces a natural $4\pi$ periodicity in the internal dynamics, providing a realist account of the spin-$\frac{1}{2}$ behavior that is traditionally ascribed to the double-valued nature of spinor representations.

The presence of this periodic internal degree of freedom, characterized by the phase \(\omega t\), defines an autonomous internal clock. This allows for a reinterpretation of gauge freedom not as redundancy in the mathematical description, but as a manifestation of the relativity of internal phase among discrete, mutually interacting quantum objects. Local $U(1)$ symmetry thus emerges from the geometric requirement of synchronizing these internal clocks, offering a realist alternative to field-theoretic formulations.

The framework makes several testable predictions that distinguish it from conventional interpretations. The Zitterbewegung velocity $v_{\text{ZB}} = 0.040374c$ should be measurable as the characteristic oscillation speed of internal phase dynamics, representing a concrete signature of the electron's internal structure. The spatial scale over which phase synchronization occurs is expected to be related to the Compton wavelength modified by the Zitterbewegung factor, providing a geometric connection between quantum mechanics and the proposed internal dynamics. Additionally, the observed deviation of the electron's magnetic moment from the classical value $g = 2$ emerges from Lorentz contraction effects in the internal oscillations, offering a particle-based explanation for this fundamental quantum correction without invoking virtual particle loops or radiative corrections in the traditional sense.



% ========================================
% Section 4: Realist Gauge Geometry from First Principles
% ========================================
\section{Realist Gauge Geometry from First Principles}

% ----------------------------------------
% Subsection 4.1: Reinterpreting U(1) Gauge Symmetry
% ----------------------------------------
\subsection{Reinterpreting $U(1)$ Gauge Symmetry}

In conventional quantum field theory, $U(1)$ gauge symmetry is introduced by promoting the global phase invariance of the wavefunction $\psi \to e^{i\alpha} \psi$ to a local symmetry $\psi(x) \to e^{i\alpha(x)} \psi(x)$. To maintain invariance under this local transformation, a gauge field $A_\mu(x)$ is introduced, transforming as $A_\mu \to A_\mu - \partial_\mu \alpha(x)$. The field strength $F_{\mu\nu} = \partial_\mu A_\nu - \partial_\nu A_\mu$ emerges as a derived quantity, and the electromagnetic interaction is described by the coupling $j^\mu A_\mu$.

However, as noted by Zee and others, this structure introduces transformations that connect not physically distinct states, but alternative mathematical descriptions of the same configuration. This has long posed a conceptual tension regarding the nature of gauge symmetry—whether it should be viewed merely as a descriptive artifact, or as a reflection of deeper physical principles. The present work seeks to resolve this ambiguity by grounding gauge structure in the observable dynamics of internal temporal phase, thereby providing a realist interpretation of gauge freedom.

In the 0-Sphere model, the structure of the theory favors a realist interpretation. Each particle carries an internal time phase $\theta(t)$, derived from a physical harmonic oscillator governed by its own Hamiltonian. The phase evolution is not a redundancy but a physical observable. Consequently, allowing $\theta$ to vary locally across an ensemble of particles introduces the need for a compensating connection—a structure that ensures phase coherence across the system. This connection is identified with the electromagnetic gauge potential $A_\mu$.


Rather than being introduced axiomatically, the gauge potential in our model emerges as a synchronization field among many internally evolving phase oscillators. The familiar $U(1)$ transformation,

\begin{equation}
\theta(x) \to \theta(x) + \alpha(x),
\end{equation}
then becomes a real physical operation: it corresponds to resetting the internal phase of a particle at spacetime point $x$.



% ----------------------------------------
% Subsection 4.2: Geometric Foundation: From Snell's Law to Geodesic Motion
% ----------------------------------------
\subsection{Geometric Foundation: From Snell's Law to Geodesic Motion}

In classical field theory, the gauge potential $A_\mu(x)$ is introduced axiomatically to maintain local phase invariance. However, this construct is inherently Eulerian—it assumes a continuous field defined over spacetime. In contrast, our approach follows the Lagrangian method of fluid dynamics: rather than defining fields at fixed spacetime points, we track individual particles as they move through spacetime, each carrying an internal phase $\theta(t)$ that evolves according to its own Hamiltonian as described by the 0-Sphere model~\cite{hanamura2024quantum,hanamura2023}. This motivates the central inquiry of this section: whether the gauge structure can emerge from the physical geometry of internal phase transport, independent of a continuous background field.

A key insight of the 0-Sphere electron model is that the transport of energy between thermal potential energy (TPE) kernels obeys a form of Snell’s law, arising not in physical space but in the internal phase space of the electron~\cite{hanamura2024quantum}. In this framework, each TPE kernel—such as Kernel $A$ located at position $+a$, and Kernel $B$ at $-a$—represents a localized bundle of mass energy that has been thermalized. These kernels act as sources and sinks of radiation-like energy that traverses a connecting structure referred to as the “photon sphere” ($\gamma^*_{\mathrm{K.E.}}$), which mediates the oscillatory energy exchange.

This internal radiative process can be described schematically as:
\begin{equation}
T_{\mathrm{kernel}A} \rightarrow \gamma^*_{\mathrm{K.E.}} \rightarrow T_{\mathrm{kernel}B},
\end{equation}
where the directionality corresponds to energy flow from the radiative source to the absorber. The resulting energy transport follows a refraction-like rule:

\begin{equation}
\frac{\sin \theta_A}{\sin \theta_B} = \frac{v_A}{v_B} \equiv n_{A \to B},
\end{equation}
where Kernels $A$ and $B$ serve as the fixed starting and ending points respectively, as determined by the principle of least action. The angles $\theta_A$ and $\theta_B$ correspond to the propagation angles of the photon sphere ($\gamma^*_{\mathrm{K.E.}}$) as it mediates radiative energy transport between these fixed endpoints. The phase velocities $v_A$ and $v_B$ characterize the local oscillatory media through which the photon sphere propagates in the vicinity of each kernel.

In this context, the media are not conventional optical substances but instead represent regions of oscillatory potential shaped by the local properties of the TPE kernels. The photon sphere follows a geodesic path between the fixed kernels $A$ and $B$, analogous to how light follows the shortest optical path in accordance with Snell's law. This geodesic motion represents the path that minimizes the action integral, embodying the principle of least action applied to internal radiative energy transfer.

Specifically, each region's phase velocity $v$ is determined by the characteristic frequency of the internal oscillations and the thermal density of the kernel. A higher local TPE concentration—i.e., a steeper thermal gradient—yields a higher-frequency Zitterbewegung component, and thus a faster internal phase velocity. The ratio $v_A / v_B$ quantifies how the wavefront of internal energy, carried by the photon sphere, refracts as it propagates through these dynamically different regions.

This interpretation gives geometric and energetic meaning to the refractive index $n_{A \to B}$, which reflects how the photon sphere—mediating internal energy transport—follows a geodesic path under the influence of asymmetric thermal distributions between the fixed kernel endpoints. The application of Snell's law in this setting is therefore not metaphorical but represents the physical principle that radiative energy follows the path of stationary action, grounded in the internal phase dynamics and thermal gradients of the electron.

This adherence to Snell's law leads to the principle of least action. The energy transport follows the path that minimizes the action integral:
\begin{equation}
S = \int_A^B ds,
\end{equation}
subject to the boundary conditions that TPE is localized at fixed points $A$ and $B$. As demonstrated by Hanamura~\cite{hanamura2024quantum}, this constraint leads directly to geodesic motion in the internal phase space, connecting quantum oscillations with the geometric principles of general relativity.




% ----------------------------------------
% Figure 1: f-TotalHamiltonian.png -- Energy conservation and hierarchical phase
% ----------------------------------------
\begin{figure*}[t!]
    \centering 
    \includegraphics[width=\textwidth]{f-TotalHamiltonian.png}
    \caption{Energy conservation and hierarchical phase structure in the 0-Sphere electron model. The graph illustrates the coexistence of $4\pi$ periodicity (fermionic kernels: blue dashed and orange dotted lines) and $2\pi$ periodicity (bosonic photon sphere: green dotted line) within a unified temporal framework. The constant total energy (red solid line, $H(\phi)=1$) demonstrates perfect energy conservation throughout the oscillation cycle. This visualization reveals how spin (half-frequency) and gauge (full-frequency) degrees of freedom emerge from a single internal clock, providing a geometric foundation for unifying fermionic and bosonic behaviors within a $U(1)$ framework.}
    \label{f-TotalHamiltonian}
\end{figure*}


% ----------------------------------------
% Subsection 4.3: Discrete Bundle Structure vs. Continuous Field Theory
% ----------------------------------------
\subsection{Discrete Bundle Structure vs. Continuous Field Theory}

The fiber bundle structure in our framework mirrors the geometric formulation of classical mechanics, where each particle carries its own tangent bundle. In classical mechanics, a point particle's motion is completely determined by its position and momentum as functions of time. When the forces acting on the particle are known, the equations of motion uniquely determine the trajectory—that is, the position and momentum for all times. This creates a curve in the space of position and momentum with time as a parameter. Geometrically, this corresponds to assigning a single vector bundle (the tangent bundle) to each particle, encoding its degrees of freedom along its worldline.

The 0-Sphere model adopts precisely this particle-centric approach. Just as classical mechanics assigns a tangent bundle to each particle to encode its kinematic degrees of freedom, our model assigns a discrete $U(1)$ fiber bundle to each individual electron to encode its internal phase degree of freedom. This bundle is attached to the particle's worldline trajectory and evolves along with the particle's motion through spacetime. This stands in fundamental contrast to conventional quantum field theory, where the $U(1)$ bundle is defined continuously over the entire spacetime manifold. In standard gauge field theory, the gauge field $A_\mu(x)$ represents a connection that exists at every point in space, creating a continuous structure that accommodates the superposition principle and field quantization procedures leading to quantum electrodynamics.

Crucially, the discrete bundles in our model do not achieve the dense coverage characteristic of quantum fields. Unlike continuous field theories where field values are defined at every spacetime point, our discrete fiber bundles exist only along particle worldlines. \textbf{The gauge connection $A_\mu$ emerges not as a fundamental field permeating spacetime, but as a geometric synchronization mechanism operating between spatially separated, discrete phase-bearing particles}—analogous to how classical forces coordinate the motion of discrete particles without requiring a continuous medium.

While our model assigns discrete internal bundle structures to individual electrons, the background spacetime itself remains a smooth, differentiable manifold. This continuity is essential for the validity of general relativity, which requires a well-defined energy-momentum tensor and the capacity to express spacetime curvature as a smooth deformation. Our proposal does not abandon this continuum, but rather considers it as the differentiable background across which discrete electrons are distributed---each acting as a localized oscillator that exchanges thermal potential and kinetic energy along a photon-mediated internal trajectory governed by a sinusoidal gradient. These internal phase dynamics, while discrete in origin, interact through the manifold by synchronizing their phases, giving rise to emergent gauge connections that geometrize inter-particle correlations without invoking a globally defined quantum field.

A differentiable manifold arises within each electron from the continuous trajectory of the photon sphere oscillating between Kernel $A$ and Kernel $B$. This internal motion, governed by a thermal gradient that follows a $\sin\theta$ profile \cite{hanamura2018}, creates a well-defined geodesic structure along which internal gauge connections can be assigned. The photon sphere does not move uniformly, but rather undergoes harmonic oscillation as it experiences varying refractive conditions along the gradient, analogous to Snell's law. A local vector bundle structure can be defined along the centerline of this trajectory, establishing a continuous geometric substrate on which internal gauge phases evolve.

The reconciliation of discrete internal structure with the smooth geometry of spacetime represents a key conceptual innovation of the present work. It suggests a realist reinterpretation in which electromagnetic fields emerge not from continuous quantum fields defined over spacetime, but from geometric relations among discrete phase-bearing particles embedded in a continuous manifold. This framework offers a novel basis for unifying the gauge principle with general relativity without appealing to field quantization, instead grounding interactions in the geometric and thermodynamic properties of structured quantum matter.


% ----------------------------------------
% Subsection 4.4: Unification without Non-Abelian Complexity
% ----------------------------------------
\subsection{Unification without Non-Abelian Complexity}

A natural objection to any gauge-theoretic unification based solely on $U(1)$ is that it cannot reproduce the rich non-Abelian structure required by the standard model, particularly for the weak and strong interactions. In conventional field theory, non-Abelian gauge groups such as $SU(2)$ and $SU(3)$ are indispensable because they allow for gauge bosons that interact with one another, producing nonlinear dynamics and self-coupling effects.

However, in the present realist framework, the role of the gauge field is reinterpreted not as a fundamental degree of freedom but as a geometric synchronization mechanism among internal temporal phases. From this perspective, interaction is not mediated by field quanta with nontrivial group structure, but rather emerges from the exchange and coordination of internal phase dynamics.

This shift allows for an alternative route to describing non-Abelian-like phenomena. For example, \textbf{internal phase spaces may be extended to include multiple independent oscillatory modes $\theta_1(t), \theta_2(t), \dots$, each corresponding to a distinct synchronization channel.} Asymmetries in the coordination of these modes, or order-dependent phase relationships, may give rise to effectively non-Abelian behavior—analogous to the commutation structure in groups like $SU(2)$—even if the underlying gauge group remains $U(1)$. Similarly, the directionality or sequential nature of energy transfer processes may encode interaction asymmetries that mimic gauge boson self-coupling in conventional theories.

Thus, while the present model does not yet explicitly incorporate non-Abelian group structures, it opens a pathway for realizing comparable dynamics through internal geometric degrees of freedom. In this view, non-Abelian gauge behavior could arise not from algebraic constraints imposed at the field level, but from emergent properties of internal temporal geometry.




% ----------------------------------------
% Table 1: Conventional Gauge Theory vs Present Realist Model
% ----------------------------------------
\begin{table*}[t!]
\renewcommand{\arraystretch}{1.5} % Adds vertical padding
\centering
\caption{Conceptual comparison between the conventional interpretation of gauge theory and the present realist model based on internal phase dynamics.}
\begin{tabular}{|p{4.5cm}|p{5.6cm}|p{6.6cm}|}
\hline
\textbf{Concept} & \textbf{Conventional Interpretation} & \textbf{Present Model (0-Sphere)} \\
\hline
Nature of electromagnetic field & Fundamental continuous field on spacetime & Geometric manifestation of internal phase synchronization among discrete particles \\
\hline
Meaning of gauge symmetry & Mathematical redundancy in field description & Physical degree of freedom associated with internal temporal phase \\
\hline
Fiber bundle structure & Continuous $U(1)$ bundle over a differentiable manifold & Discrete $U(1)$ fiber assigned to each internally oscillating electron \\
\hline
Gauge potential $A_\mu$ & Introduced to maintain local phase invariance & Connection enforcing phase coherence across distributed internal clocks \\
\hline
Field strength $F_{\mu\nu}$ & Curvature of the gauge connection & Holonomy induced by desynchronization of internal phases along closed paths \\
\hline
\end{tabular}
\label{tab:conceptual_comparison}
\end{table*}



% ----------------------------------------
% Subsection 4.5: Hierarchical Phase Periodicity as a Bridge Between Spin and Gauge
% ----------------------------------------
\subsection{Hierarchical Phase Periodicity as a Bridge Between Spin and Gauge}

An essential structural feature of the present model lies in the coexistence of two distinct temporal periodicities embedded within a single internal energy identity. As expressed in Eq.~\eqref{eq:energy_conservation}, the internal dynamics of the electron involve oscillatory components at both $\omega t$ and $\omega t/2$:

\begin{equation}
E_0 = E_0\left[\cos^4\left(\frac{\omega t}{2}\right) + \sin^4\left(\frac{\omega t}{2}\right) + \frac{1}{2}\sin^2(\omega t)\right].
\label{eq:energy_conservation}
\end{equation}

The energy identity of Eq.~\eqref{eq:energy_conservation} and its constituent oscillatory terms are visualized in Fig.~\ref{f-TotalHamiltonian}, where the $4\pi$ periodicity of the kernel components and the $2\pi$ periodicity of the radiative photon sphere are depicted as phase-dependent energy distributions. This figure supports the interpretation that the spin and gauge components, though differing in periodicity, emerge from a unified temporal evolution—a physical basis for treating them as harmonic layers of the same internal geometry.

Here, the $\omega t/2$ terms are associated with the internal dynamics of the fermionic kernel and give rise to $4\pi$ periodicity, while the $\omega t$ term corresponds to the photon sphere—the radiative structure surrounding the kernel—and exhibits $2\pi$ periodicity. Remarkably, both components evolve under the same temporal parameter $t$, governed by a common internal clock. This implies that the fermionic and bosonic characters of the electron are not fundamentally disjoint, but rather represent harmonically related expressions of a unified temporal phase geometry. Such a structure allows the spin and gauge degrees of freedom to be viewed not as separate fields or symmetries, but as multiple scales of a single phase-coherent system.

This perspective gains further support from a reinterpretation of the origin of spin angular momentum. In relativistic kinematics, an electron undergoing acceleration experiences a precession of its local frame relative to the laboratory frame. This effect, known as Thomas precession~\cite{Thomas1926, thomas1927}, leads to an angular velocity of the internal coordinate system given by~\cite{hanamura2023}
\begin{equation}
\boldsymbol{\Omega} = \frac{1}{2c^2}[\boldsymbol{a} \times \boldsymbol{v}],
\label{eq:Thomas}
\end{equation}
where $\boldsymbol{a}$ is the acceleration and $\boldsymbol{v}$ is the velocity of the electron. In the 0-Sphere model, the internal kernel undergoes simple harmonic oscillation (not uniform circular motion) with $\theta(t) = \omega t$. Substituting the corresponding acceleration and velocity of this harmonic oscillator into Eq.~\eqref{eq:Thomas} yields a remarkable result: the precessional angular velocity contains a $\sin(2\omega t)$ term, exhibiting precisely double the frequency of the underlying oscillation. This frequency doubling provides a natural geometric explanation for spin-1/2 quantization—the requirement of a 720° rotation for the electron to return to its original state arises not from abstract quantum postulates, but from the relativistic kinematics of the internal structure~\cite{hanamura2023}.

Hence, the spin emerges not as a static intrinsic quantum number, but as a dynamic consequence of precession-driven phase accumulation at $\omega t/2$, consistent with the $4\pi$ periodicity of fermionic behavior.

This construction reveals that the same internal phase variable $\theta(t)$, when expressed at different harmonic levels, accounts simultaneously for the spin degree of freedom and the electromagnetic interaction. The gauge field $A_\mu$ arises to synchronize the $\omega t$ dynamics of the radiative shell across particles, while the spin emerges from the $\omega t/2$ component tied to kernel precession. Both are governed by the same time-dependent oscillator, leading to a hierarchy of phase structures encoded within a single $U(1)$ geometry.

The implication is that the diversity of interaction channels—traditionally modeled by multiple gauge symmetries such as $SU(2)$ and $SU(3)$—might instead arise from structured internal phase harmonics within a unified $U(1)$ framework. The nontrivial behavior typically attributed to non-Abelian symmetry groups could be emergent from the dynamics of internal phase synchronization among layered temporal modes. In this sense, what appears as group-theoretic complexity may reflect the rich structure of hierarchical phase geometry internal to each particle.



\newpage



% ========================================
% Section 5: Conclusion
% ========================================
\section{Conclusion}

This work proposes a reconceptualization of gauge theory based on internal temporal structure. Departing from the conventional view that local gauge symmetry reflects a formal redundancy in field-theoretic descriptions, the present model interprets it as a geometric necessity arising from phase synchronization among discrete, oscillating particles. The time phase $\theta(t)$ is not an auxiliary construct but a physically meaningful degree of freedom whose evolution governs electromagnetic phenomena.

In this realist framework, gauge freedom is not a redundancy to be eliminated but an expression of the intrinsic autonomy of each particle's internal temporal state. The gauge potential $A_\mu$ arises as the geometric connection required to ensure covariant transport of internal temporal phase across spacetime, within a discrete ensemble of locally oscillating particles. Rather than being imposed externally, $A_\mu$ encodes the necessary connection for synchronizing these internal time phases across spacetime.

This interpretation is summarized conceptually in Table~\ref{tab:conceptual_comparison}, which contrasts the conventional field-theoretic approach with the present particle-based model. While traditional gauge theories derive interactions from imposed symmetries on continuous fields, the present framework derives them from the coordination of internal oscillations among discrete phase-bearing particles. \textbf{By attributing the emergence of gauge degrees of freedom to physically observable internal structure---specifically, temporal phase dynamics induced by spacetime deformation---}this model offers one of the first concrete attempts to resolve the long-standing problem of gauge redundancy through physical geometry.

Furthermore, the inclusion of internal harmonic motion within each electron allows the model to accommodate not only special relativistic effects such as Lorentz contraction, but also geometric principles from general relativity, including geodesic precession. This embedding of both relativistic frameworks into a single particle-based structure underscores the foundational scope of the internal phase geometry proposed here.

Of particular note, the model offers an experimentally testable prediction: the internal Zitterbewegung velocity $v_{\text{ZB}} = 0.040374c$, obtained from a closed-form algebraic framework grounded in special relativity~\cite{hanamura2025bridging}. This value, though not yet experimentally verified, is a first-principles calculation refined by incorporating the concept of geodesic precession from general relativity. This synthesis demonstrates how quantum behavior can emerge from deterministic internal dynamics subject to relativistic geometric constraints, offering a new algebraic bridge between quantum mechanics and gravity—beyond the language of gauge symmetry alone.

This framework establishes the foundational principle that discrete internal structures---localized fiber bundles encoding internal phase dynamics---can coexist with continuous spacetime geometry. While each electron carries its own autonomous phase evolution, the global coordination of such discrete fibers across the manifold remains an open theoretical frontier. By recasting gauge interaction as geometric synchronization among these internal clocks, the model offers a particle-based pathway toward unifying quantum mechanics and general relativity without invoking field quantization.

As Zee has suggested~\cite{zee2010}, the redundancy inherent in current gauge theory formulations may eventually give way to a more physically grounded approach. The present work aims to advance that transition, showing that local symmetry may emerge not from formal constraints, but from the geometry of internal phase evolution itself.


\newpage



% Replace existing thebibliography with inline entries (references.bib not used)
%\bibliographystyle{apsrev4-1}
%  APS Physical Review style
%\bibliography{references}
%  Would reference references.bib file

% ========================================
% Bibliography
% ========================================
\begin{thebibliography}{99}
\bibitem{zee2010} Zee, A., 
\textit{Quantum Field Theory in a Nutshell}, 
2nd ed., Princeton University Press, Princeton (2010), p. 474. 
ISBN: 978-0-691-14034-6

\bibitem{wilson1974} Wilson, K. G., 
\textit{Confinement of quarks}, 
Phys. Rev. D \textbf{10}, 2445--2459 (1974). 
\href{https://link.aps.org/doi/10.1103/PhysRevD.10.2445}{doi:10.1103/PhysRevD.10.2445}.

\bibitem{rovelli1995} Rovelli, C. and Smolin, L., 
\textit{Spin networks and quantum gravity}, 
Phys. Rev. D \textbf{52}, 5743--5759 (1995). 
\href{https://arxiv.org/abs/gr-qc/9505006}{arXiv:gr-qc/9505006}.

\bibitem{hanamura2024quantum} 
S. Hanamura, Quantum Oscillations in the 0-Sphere Model: Bridging Zitterbewegung, Geodesic Paths, and Proper Time Through Radiative Energy Transfer, Zenodo (2024), \href{https://doi.org/10.5281/zenodo.17765017}{doi:10.5281/zenodo.17765017}. [Originally: \href{https://vixra.org/abs/2411.0117}{viXra:2411.0117}]

\bibitem{aharonov1959} Aharonov, Y. and Bohm, D.,
\textit{Significance of electromagnetic potentials in the quantum theory},
Phys. Rev. \textbf{115}, 485--491 (1959).

\bibitem{weyl1929} Weyl, H.,
\textit{Elektron und Gravitation. I},
Z. Phys. \textbf{56}, 330--352 (1929).

\bibitem{yang1954} Yang, C. N. and Mills, R. L.,
\textit{Conservation of isotopic spin and isotopic gauge invariance},
Phys. Rev. \textbf{96}, 191--195 (1954).

\bibitem{tonomura1986} Tonomura, A., Osakabe, N., Matsuda, T., Kawasaki, T., Endo, J., Yano, S., and Yamada, H.,
\textit{Evidence for Aharonov-Bohm effect with magnetic field completely shielded from electron wave},
Phys. Rev. Lett. \textbf{56}, 792--795 (1986).

\bibitem{teller2000} Teller, P.,
\textit{The gauge argument},
Philosophy of Science \textbf{67}, S466--S481 (2000).
\href{https://doi.org/10.1086/392824}{doi:10.1086/392824}

\bibitem{hanamura2018} 
S. Hanamura, A Model of an Electron Including Two Perfect Black Bodies, Zenodo (2018), \href{https://doi.org/10.5281/zenodo.16759284}{doi:10.5281/zenodo.16759284}. [Originally: \href{https://vixra.org/abs/1811.0312}{viXra:1811.0312}]

\bibitem{berry1984} Berry, M. V.,
\textit{Quantal phase factors accompanying adiabatic changes},
Proc. R. Soc. Lond. A \textbf{392}, 45--57 (1984).
\href{https://doi.org/10.1098/rspa.1984.0023}{doi:10.1098/rspa.1984.0023}

\bibitem{zitterbewegung_schrodinger1930} Schrödinger, E.,
\textit{Über die kräftefreie Bewegung in der relativistischen Quantenmechanik},
Sitzungsber. Preuss. Akad. Wiss. Phys.-Math. Kl. \textbf{24}, 418--428 (1930).

\bibitem{hanamura2023} 
S. Hanamura, Redefining Electron Spin and Anomalous Magnetic Moment Through Harmonic Oscillation and Lorentz Contraction, Zenodo (2023), \href{https://doi.org/10.5281/zenodo.17764997}{doi:10.5281/zenodo.17764997}. [Originally: \href{https://vixra.org/abs/2309.0047}{viXra:2309.0047}]

\bibitem{Thomas1926} Thomas, L. H., 
\textit{The Motion of the Spinning Electron}, 
Nature \textbf{117}, 514 (1926). DOI: \href{https://doi.org/10.1038/117514a0}{10.1038/117514a0}

\bibitem{thomas1927} Thomas, L. H., 
\textit{The Kinematics of an Electron with an Axis}, 
Philosophical Magazine \textbf{3}, 1 (1927). DOI: \href{https://doi.org/10.1080/14786440108564170}{10.1080/14786440108564170}

\bibitem{hanamura2025bridging} 
S. Hanamura, Bridging Quantum Mechanics and General Relativity: A First-Principles Approach to Anomalous Magnetic Moments and Geodetic Precession, Zenodo (2025), \href{https://doi.org/10.5281/zenodo.17765097}{doi:10.5281/zenodo.17765097}. [Originally: \href{https://vixra.org/abs/2501.0130}{viXra:2501.0130}]

\end{thebibliography}



% = = * * = =  * * = = * * = = * * = = * * = = * * = = * * = = * * 
% \section{\label{sec:Appendix}Appendix}

% \section*{Acknowledgments}

\end{document}